\documentclass[conference]{IEEEtran}

\usepackage{cite}
\usepackage{amsmath,amssymb,amsfonts}
\usepackage{algorithm}
\usepackage{algorithmic}
\usepackage{graphicx}
\usepackage{textcomp}
\usepackage{booktabs}
\usepackage[capitalize,noabbrev]{cleveref}

\setlength{\textfloatsep}{4pt plus 1pt minus 1pt}
\setlength{\floatsep}{4pt plus 1pt minus 1pt}
\setlength{\intextsep}{4pt plus 1pt minus 1pt}
\setlength{\abovedisplayskip}{4pt plus 1pt minus 1pt}
\setlength{\belowdisplayskip}{4pt plus 1pt minus 1pt}
\setlength{\abovedisplayshortskip}{2pt plus 1pt minus 1pt}
\setlength{\belowdisplayshortskip}{2pt plus 1pt minus 1pt}

\IEEEoverridecommandlockouts

\begin{document}

\title{Optimal Carbon Pricing of Closed-Loop Supply Chains: A
Global-Optimisation of a Bi-Level Mixed-Integer Problem
\thanks{Corresponding author: M. Pinto (marcelo.pinto@dem.uc.pt).}}

% IEEEtran's \and glues independent, content-width boxes with stretchable
% space rather than laying out a real grid, so uneven affiliation lengths
% make the rows drift out of alignment. A nested tabular cannot fix this
% either: \@author is expanded directly inside IEEEtran's own internal
% \halign, and a second \halign/tabular nested there corrupts that
% alignment outright. The robust fix is to avoid any nested alignment and
% build each cell as a fixed-width \parbox, paired left/right with
% \makebox, so column edges are pinned by box width rather than content.
\newcommand{\authorcell}[2]{\makebox[0.47\textwidth]{\parbox[t]{0.44\textwidth}{\centering{\sublargesize #1}\\{\normalsize #2}}}}

\author{
\authorcell{Marcelo Pinto}{CEMMPRE, ARISE, Dept.\ Mech.\ Eng., Univ.\ Coimbra\\Coimbra, Portugal}\hfill
\authorcell{Paulo Nascimento}{CEMMPRE, ARISE, Dept.\ Mech.\ Eng., Univ.\ Coimbra\\Coimbra, Portugal}\\[10pt]
\authorcell{Crist\'ov\~ao Silva}{CEMMPRE, ARISE, Dept.\ Mech.\ Eng., Univ.\ Coimbra\\Coimbra, Portugal}\hfill
\authorcell{Samuel Moniz}{CEMMPRE, ARISE, Dept.\ Mech.\ Eng., Univ.\ Coimbra\\Coimbra, Portugal}
}

\maketitle

\begin{abstract}
We study the design of closed-loop supply chains under a
regulator-imposed carbon tax, formulated as a bilevel mixed-integer
program~(BL-MIP). The regulator~(leader) selects a discrete
carbon-tax rate that minimises supply-chain emissions; the
operator~(follower) responds by optimising facility location,
capacity, technology, transportation, and forward and reverse product
flows to minimise cost inclusive of tax. We solve the resulting
BL-MIP with a deterministic bounding procedure~(DBP) that computes
the regulator's optimum between two bounds: a linear-programming~(LP)
relaxation of the follower's problem provides a lower bound, and the 
follower's exact integer response provides an upper bound. The DBP
iteratively tightens these two bounds until their gap is below a 
given tolerance. Three results emerge. First, emission abatement 
saturates at a network-specific threshold, beyond which higher taxes 
yield no further reduction. Second, abatement is governed by the
technology menu --- cleaner production and remanufacturing options
--- and expanding that menu can reduce both cost and emissions. 
Finally, expanding recovery capacity achieves the same abatement 
at a lower carbon price --- each unit of recovery substitutes 
directly for a unit of taxation.
\end{abstract}

\begin{IEEEkeywords}
bilevel optimization, carbon pricing, closed-loop supply chain network
design, mixed-integer programming, environmental regulation
\end{IEEEkeywords}

\section{Introduction}
\label{sec:intro}

Carbon taxation prices each unit of CO\textsubscript{2}-equivalent
emitted, producing emission cuts when well calibrated \cite{pigou1920}, 
and closed-loop architectures that integrate recovery and remanufacturing 
enlarge the available abatement levers \cite{govindan2015}. Crucially, 
the regulator-operator interaction is hierarchical: the regulator 
commits to a tax before the operator responds, so a single-level model
that treats the tax as fixed cannot represent the regulator's problem of
choosing a rate that induces a desirable response.

We model this as a bilevel program \cite{colson2007}: the leader sets
the carbon tax to minimize total emissions and, among rates achieving
that minimum, levies the lowest one; the follower solves a
multi-echelon, multi-period design problem with integer facility,
technology, and transport decisions, minimizing cost inclusive of tax.
The resulting bilevel mixed-integer program~(BL-MIP) is NP-hard 
\cite{jeroslow1985}.

This paper makes three contributions: (i) a BL-MIP with a discrete 
carbon tax decision with a closed-loop supply-chain design problem;
(ii) a deterministic bounding procedure (DBP) converging to a 
$\epsilon$-optimal solution; and (iii) computational results show 
that emission abatement saturates at a network-specific tax threshold,
that technology is the binding lever below that threshold, and that 
expanding recovery capacity can substitute directly for taxation. 
\Cref{sec:model} states the model, \Cref{sec:method} the DBP, and 
\Cref{sec:results} the results.

\section{Bilevel Model}
\label{sec:model}

The closed-loop network includes suppliers, factories, warehouses,
air/sea hub terminals, and markets, linked by forward flows of raw
materials and finished products and by reverse flows of end-of-life
items recovered for remanufacturing. Over a multi-period planning
horizon, the follower decides the network structure --- facility, 
technology, transport, and fleet choices --- and production,
remanufacturing, inventory, and flow quantities. 

Let $\mathbf{z}\in\mathcal{X}$ denote the follower's decision 
vector, where $\mathcal{X}$ is the follower's feasible region 
induced by (i) flow and inventory balance, (ii) reverse-logistics 
and returns, (iii) raw-material supply, facility-opening, (iv) 
area-, inventory and fleet-capacity, (v) technology-, and 
transport-selection, and (vi) intermodal hub routing. The 
regulator's decision variable is the carbon-tax rate 
$\tau$ (USD/tCO\textsubscript{2}e), subject to integer values in
$\mathcal{T}=\{\tau_{\min},\tau_{\min}+1,\ldots,\tau_{\max}\}$. 
Denoting the follower's cost and emissions by $\mathrm{Cost}(\mathbf{z})$
and $\mathrm{Emis}(\mathbf{z})$, its optimal response to a given rate
$\tau$ is the correspondence
\begin{equation}
\mathcal{S}(\tau) \;\triangleq\;
\operatorname*{arg\,min}_{\mathbf{z}\in\mathcal{X}}\;
\bigl\{\mathrm{Cost}(\mathbf{z}) + \tau\,\mathrm{Emis}(\mathbf{z})\bigr\}.
\label{eq:follower}
\end{equation}

Anticipating this response, the leader minimizes resulting emissions and, 
among rates attaining that minimum, the tax itself (avoiding unnecessary 
taxation). This defines the bilevel program
\begin{subequations}
\label{eq:leader}
\begin{alignat}{2}
\textup{(BL)}\quad
&\operatorname*{lexmin}_{\tau,\,\mathbf{z}}\quad
  \bigl(\mathrm{Emis}(\mathbf{z}),\ \tau\bigr) \\
&\text{s.t.}\quad \tau\in\mathcal{T}, \qquad
  \mathbf{z}\in\mathcal{S}(\tau).
\end{alignat}
\end{subequations}
Since the follower's problem~\eqref{eq:follower} is itself a
mixed-integer program for every fixed $\tau$, problem~\eqref{eq:leader}
is at least as hard, establishing the NP-hardness of the BL-MIP noted
in \Cref{sec:intro} \cite{jeroslow1985}.

\section{Solution Method}
\label{sec:method}

We solve~\eqref{eq:leader} with a DBP, adapted from \cite{soares2021}, 
which generates avsequence of converging lower and upper bounds on the 
leader's optimalvemissions, in the following steps.

\textbf{Step 1} (initialization). Set $\mathrm{LB}\leftarrow-\infty$.
Solve the \emph{relaxation} of~\eqref{eq:leader} obtained by dropping
the coupling constraint $\mathbf{z}\in\mathcal{S}(\tau)$, i.e.,
$\min\{\mathrm{Emis}(\mathbf{z}):\mathbf{z}\in\mathcal{X},\,\tau\in\mathcal{T}\}$,
to obtain a first candidate rate $\tau$.

\textbf{Step 2} (follower response). Solve the follower's exact
problem~\eqref{eq:follower} at $\tau$ to obtain its true optimum
$\nu(\tau)$.

\textbf{Step 3} (optimistic tie-break). Since $\mathcal{S}(\tau)$
may contain several optimal solutions, solve
$\min\{\mathrm{Emis}(\mathbf{z}):\mathbf{z}\in\mathcal{X},\,
\mathrm{Cost}(\mathbf{z})+\tau\,\mathrm{Emis}(\mathbf{z})\le\nu(\tau)\}$
to recover its lowest-emissions member $\mathbf{z}^{*}$.

\textbf{Step 4} (upper bound). The pair $(\tau,\mathbf{z}^{*})$ is
feasible for~\eqref{eq:leader}. If $\mathrm{Emis}(\mathbf{z}^{*})$
improves on the incumbent, update $\mathrm{UB}\leftarrow
\mathrm{Emis}(\mathbf{z}^{*})$ and $\tau^{\dagger}\leftarrow\tau$.

\textbf{Step 5} (tightening). Add to the relaxation of Step~1 the cut
$\mathrm{Cost}(\mathbf{z})+\tau\,\mathrm{Emis}(\mathbf{z})\le
\mathrm{Cost}(\mathbf{z}^{*})+\tau\,\mathrm{Emis}(\mathbf{z}^{*})$
and the bound $\mathrm{Emis}(\mathbf{z})\le\mathrm{UB}$, ruling out
$\tau$ from being proposed again; re-solve it for a new candidate
$\tau$ and update $\mathrm{LB}$.

\textbf{Step 6} (stopping test). Stop and return $\tau^{\dagger}$ if
$(\mathrm{UB}-\mathrm{LB})/\mathrm{UB}\le\epsilon$, or if neither bound
has improved over $10$ consecutive iterations; otherwise return to
Step~2.


\section{Preliminary Results}
\label{sec:results}

We generated eight instances, following the instance-generation
guidelines of \cite{tautenhain2021}: network (17/25 nodes), product 
(4/8 types), and horizon (3/6 periods). All 8 instances converge with
a relative gap no greater than $\le1.36\times10^{-5}$ in 2--4 DBP
iterations (\Cref{tab:results}). We also run two further experiments. 
A \emph{scenario analysis} re-solves each base instance under two menu expansions 
-- technology  (3$\to$6 production and  3$\to$6 remanufacturing technologies) and 
transport (6$\to$12 modes, two additional road/air/sea options each) -- for 16 runs. A 
\emph{sensitivity analysis} then perturbs each base instance one factor at 
a time (OAT) across 11 model factors, grouped as: demand; production and 
transportation capacity; warehouse and air/sea hub cost; return availability; 
return cost; remanufacturing bill-of-materials; production and upstream supply 
emissions; and carbon-tax bounds -- for 152 runs.

\begin{table}[t]
\centering
\caption{DBP solution per base instance.}
\label{tab:results}
\renewcommand{\arraystretch}{1.0}
\footnotesize
\begin{tabular}{@{}l@{\hspace{10pt}}r@{\hspace{12pt}}r@{\hspace{6pt}}r@{}}
\toprule
\textbf{Instance} & $\tau^{\star}$ & $\Delta E$ (\%) & $\Delta C$ (\%) \\
\midrule
I17 M4 T3 & 200 & $-28.4$ & $+14.2$ \\
I17 M4 T6 & 190 & $-9.3$  & $+3.4$  \\
I17 M8 T3 & 140 & $-23.9$ & $+75.6$ \\
I17 M8 T6 & 200 & $-8.7$  & $+5.0$  \\
I25 M4 T3 & 180 & $-35.4$ & $+75.9$ \\
I25 M4 T6 & 20  & $-57.6$ & $+20.6$ \\
I25 M8 T3 & 140 & $-48.4$ & $+14.9$ \\
I25 M8 T6 & 100 & $-8.6$  & $+5.2$  \\
\bottomrule
\end{tabular}

\vspace{2pt}
\noindent\footnotesize $\tau^{\star}$ in USD/tCO\textsubscript{2}e;
$\Delta E$ and $\Delta C$ are the percentage changes in emissions
and cost, respectively, relative to each instance's $\tau=0$
baseline.
\end{table}

Four findings emerge from our runs.

\emph{First}, carbon pricing reduces emissions by $8.6$--$57.6\%$ at
a cost increase of $3.4$--$75.9\%$, and the two are not consistently
related: I25\,M4\,T6 achieves the largest cut ($-57.6\%$) at the
lowest abatement cost (\$3.62/t), whereas I17\,M8\,T3 and
I25\,M4\,T3 incur a cost increase above $75\%$ for a comparatively
modest $24$--$35\%$ reduction.

\emph{Second}, because the operator chooses among finitely many
configurations, its emission response saturates beyond a threshold:
widening the ceiling from \$200/t to \$800/t moves I25\,M8\,T3's
optimum eighteen-fold (\$20$\to$\$360/t) for a further $0.04\%$ cut,
while I17\,M8\,T3 stays flat across \$140--\$310/t -- risking
arbitrarily high taxes for no added abatement.

\emph{Third}, technology, is the binding lever:
widening the technology menu (3$\to$6/family) cuts emissions in all
8 instances ($-18.9\%$ to $-37.8\%$) and, on I25\,M4\,T3, cost by
$41.8\%$ -- a Pareto gain regardless of $\tau^{\star}$. Widening
transport modes (6$\to$12) instead leaves emissions near-unchanged
($0.24\%$).

\emph{Fourth}, recovery capacity substitutes for carbon pricing:
greater return availability lowers the tax rate needed to reach a
given emissions level (elasticity $-0.744$) without changing that
level itself (elasticity $-0.001$). More recoverable material lets
the cheaper remanufacturing pathway absorb output with less tax
pressure -- recovery substitutes for, rather than competes with,
taxation.

\section{Conclusion and Ongoing Work}
\label{sec:conclusion}

We are also solving a multi-objective formulation in which 
the carbon tax is fixed and cost is traded off against emissions, 
to quantify the gain from the bilevel program.

{\footnotesize
\bibliographystyle{IEEEtran}
\bibliography{references}
}

\end{document}
